\par\bigskip\Needspace{8\baselineskip}
\pdfbookmark[1]{Perturbations of Ceres: Remarks}{ceres-bemerkungen}
\SourceLeaf{416}
\begin{center}REMARKS.\end{center}
\begingroup\small
\noindent\hspace*{\parindent}Note [2] of section I comes from a booklet entitled Schedae, and notes [3]---[7] from a notebook. The fragmentary state of the originals required a rather free editorial treatment. The notebook also contains the first numerical calculation of the perturbations using the formulas given in [7]. Its results are printed in \textsc{Zach}'s Monatliche Correspondenz, vol. VI (Werke, vol. VI, p. 225). On the basis of these perturbations, \textsc{Gauss} again improved the elements (vol. VI, p. 228 and note [1]), and then recalculated the perturbations with the improved elements. This latter calculation is found on several leaves stitched together under the title »Praktische Anweisung zur Berechnung der Störungen«. They also contain at the beginning a collection of the formulas used, including those for the perturbations in latitude (note [8]). \textsc{Gauss} had set the latter aside in the first calculation. Their derivation could not be located. The results of the second calculation are printed in \textsc{Zach}'s Monatliche Correspondenz, vol. VI (Werke, vol. VI, pp. 228---229). \textsc{Gauss} took Jupiter's mass to be \(\dfrac{1}{1067{,}09}\) (following \textsc{Laplace}).

The passages from letters printed under II refer to the tables that \textsc{Gauss} prepared in connection with the calculations just mentioned and published in \textsc{Zach}'s Monatliche Correspondenz, vol. VII (Werke, vol. VI, pp. 235---243).

In the interval between the two letters to \textsc{Olbers} of 25 January and 10 May 1805 (or earlier), \textsc{Gauss} began to extend his first method to the second powers of the eccentricities. The beginning of these expansions is also found in the notebook already mentioned. \textsc{Gauss} soon abandoned them because of their length, and the calculation was then made by the second method presented under III. See in this connection the letter to \textsc{Olbers} of 10 May 1805.

In this second method, \textsc{Gauss} already applied the expansion of the disturbing function by interpolation later used by \textsc{Hansen} (cf. \textsc{Gauss} to \textsc{Hansen}, 11 March 1843). The method also underlies the later calculation of the perturbations of Pallas. Only the numerical calculations could be found in the same notebook, and the expansions were reconstructed from them. As notes [2]---[5] were therefore written entirely afresh, they are printed in small type. The only material from the original is the values of \(N\), \(\sin J\), \(\AscendingNode\), and the explanation of the meanings of \(\alpha,\epsilon,\epsilon'\) on p. 402, equation 2) and the value of the constant in equation 3) on p. 403, and the numerical values on pp. 405---407, with the exception of those on p. 405, lines 9---5 from the bottom, and of the coefficients of \(\cos(3D+\text{\AstroSymbols ⚳}-\AscendingNode+M)\) on p. 406, lines 11---15, which were supplied by the editor. Here \textsc{Gauss} calculated only the perturbations in latitude. Only the very beginning of the application to the perturbations of the radius vector and longitude was found, and \textsc{Gauss} then turned to the perturbations of Pallas. The results are given in the letter to \textsc{Olbers} of 2 July 1805. They did not appear in print at the time.

On the whole, \textsc{Gauss}'s work on Ceres has the character of preliminary work for his extensive investigations of Pallas, which helps explain why certain points are treated more fully there than here.

See also the individual communications on the progress of \textsc{Gauss}'s investigations already printed in vol. VI. Minor inaccuracies were corrected in several places.
\par\nobreak\smallskip
{\raggedleft\textsc{Brendel}.\par}
\nobreak\smallskip
{\centering\rule{.38\linewidth}{.4pt}\par}
\endgroup
